\chapter{Architecture}

\section{Compute: Neural Multiprocessor}
Each of the \code{N\_PROCESSORS} Neural Processors is an independent
INT8 datapath: \code{P\_IN}=8 parallel multipliers feeding a balanced
binary adder tree into a 32-bit accumulator, followed by ReLU and INT8
saturation -- byte-for-byte identical to \code{neural\_processor.v},
unmodified since before the V2 single-SDRAM freeze. Job dispatch is a
real dependency-graph pipeline:

\begin{center}
\begin{tikzpicture}[node distance=5mm and 6mm,font=\scriptsize]
  \node[fnreg] (host) {SPI host\\bridge};
  \node[fnreg,right=of host] (dm) {Dependency\\Manager};
  \node[fnreg,right=of dm] (dir) {Neural\\Director};
  \node[fnreg,below=of dir] (mm) {Memory Manager\\(per slot)};
  \node[fnreg,left=of mm] (np) {Neural\\Processor};
  \draw[fnarrow] (host) -- (dm);
  \draw[fnarrow] (dm) -- (dir);
  \draw[fnarrow] (dir) -- (mm);
  \draw[fnarrow] (mm) -- (np);
\end{tikzpicture}
\end{center}

The Dependency Manager holds a per-node table (state, required/
resolved producer count, job descriptor fields) and hands READY nodes
to the Director one at a time (valid/ready, backpressure-safe). The
Director allocates the first free processor slot (first-free, not
load-balanced) and frees it the instant that slot's job completes.

\section{Unified single-SDRAM memory}
\begin{center}
\begin{tikzpicture}[node distance=6mm and 10mm,font=\footnotesize]
  \node[fnblockT,minimum width=50mm,minimum height=20mm] (fpga){};
  \node[anchor=north,font=\bfseries,text=fnDark] at (fpga.north){FPGA};
  \node[fnreg,fill=white] (np) at ([yshift=2mm]fpga.center){4$\times$ Neural Processor};
  \node[fnreg,fill=white,below=2mm of np] (be){Unified SDRAM Backend / Arbiter};
  \node[fnblock,right=16mm of fpga,minimum height=20mm] (ram){\code{AS4C4M16SA-6TIN}\\Weights\\Activations\\Results};
  \draw[fnbus] (fpga.east|-be) -- node[fnlbl,above]{16-bit SDRAM bus} (ram);
\end{tikzpicture}
\end{center}

\code{sdram\_unified\_backend.v} owns exactly one \code{sdram\_
controller.v} instance and presents two logical ports:
\begin{itemize}
\item \textbf{W} (weight fetch): 64-bit, read-only, a 4-entry
  fully-associative ``other half'' cache (round-robin eviction --
  safe under any sizing, since an evicted-too-early entry only costs
  an extra real fetch, never wrong data).
\item \textbf{AR} (activation fill + result write-back): 16-bit,
  read/write, byte-maskable via real SDR SDRAM DQM semantics
  (\code{lb\_n}/\code{ub\_n}).
\end{itemize}
Both ports are, in turn, each arbitrated across the N\_PROCESSORS
slots by a generic, reused \code{slot\_mem\_arbiter} (a proven,
pending-latch-based, single-owner-until-ready design used identically
throughout this project).

\subsection*{Official V2 memory map}
\begin{tabularx}{\textwidth}{L{3.2cm}L{3.2cm}Y}
\toprule
\rowh \thd{Region} & \thd{Base address} & \thd{Notes} \\
\midrule
Weights & \code{0x010000} & 1\,MB-aligned \\
\rowa Activations & \code{0x200000} & 1\,MB-aligned \\
Results & \code{0x300000} & 1\,MB-aligned \\
\bottomrule
\end{tabularx}
All three regions are non-overlapping within the single 8\,MB
(\code{0x000000}--\code{0x7FFFFF}) SDRAM address space. Base addresses
are host-programmable per job (via \code{WRITE\_JOB}'s own
\code{x\_base}/\code{w\_base}/\code{result\_addr} fields), not hard-coded
in the datapath.

\section{Read-timing discipline (ERR-0025)}
The shared weight and activation SRAMs (\code{nms\_weight\_packed.v},
\code{nms\_activation\_replicated.v}) use \textbf{combinational} read
ports, matching the exact 1-cycle latency assumption of the per-slot
Memory Manager's own pipelined read-ahead consumer. An earlier,
registered-read implementation added one uncounted cycle of latency
that a busy, multi-tile job's own prefetch lead time always absorbed
invisibly, but that an uncontested single-tile job exposed on its
first (only) tile -- found and fixed via this revision's own
board-level SPI integration testing, with zero regression to the
existing bit-exact regression suite (identical cycle counts before and
after the fix).
